%% kapa: exam=1 round=37 year=2026 period=2 form=A subject=appraisal_law subject_ko=감정평가관계법규 seq=20 number=20
%% source: Q-Net 감정평가사 기출문제 (cst00309) artlSeq=5258924 file=raw/1008/5258924/2026년 제37회 감정평가사 1차 2교시.pdf
%% conversion: kordoc 4.24.51 -> tools/gen1.py (auto); review=proofread-sonnet; flags=none
\begin{question}{20}
국유재산법령상 수의의 방법으로 행정재산의 사용허가를 받을 자를 결정할 수 있는 경우를 모두 고른 것은? (단, 주어진 조건 외에 다른 조건은 고려하지 않음)
\begin{framedtext}
ㄱ. 경작용으로 실경작자에게 사용허가를 하는 경우 \par
ㄴ. 국가와 재산을 공유하는 자에게 국가의 지분에 해당하는 부분에 대하여 사용허가를 하는 경우 \par
ㄷ. 토지의 용도 등을 고려할 때 해당 재산에 인접한 토지의 소유자를 지명할 필요가 있는 경우 \par
ㄹ. 국유재산의 관리ㆍ처분에 지장이 없는 경우로서 사용목적이나 계절적 요인 등을 고려하여 1년의 사용허가를 하는 경우 \par
\end{framedtext}
\begin{choices}
  \item ㄱ
  \item ㄱ, ㄴ
  \item ㄷ, ㄹ
  \item ㄴ, ㄷ, ㄹ
  \item ㄱ, ㄴ, ㄷ, ㄹ
\end{choices}
\end{question}
\answer{2}
